\DocumentMetadata{
  pdfversion=2.0,
  pdfstandard=ua-2,
  lang=en-US,
  tagging=on,
  tagging-setup={math/setup=mathml-SE,tabsorder=structure}
}
\documentclass[11pt,letterpaper]{article}
\usepackage[margin=0.68in,includefoot]{geometry}
\usepackage{newcomputermodern}
\usepackage{amsmath,amscd,mathtools}
\providecommand{\square}{\mdlgwhtsquare}
\usepackage[hidelinks]{hyperref}
\hypersetup{
  pdftitle={Combinatorics PhD Exam, May 1995},
  pdfauthor={University of Florida Department of Mathematics},
  pdfsubject={Graduate mathematics examination},
  pdfdisplaydoctitle=true,
  bookmarksopen=true
}
\setcounter{secnumdepth}{0}
\setlength{\parindent}{0pt}
\setlength{\parskip}{0.28em}
\setlength{\emergencystretch}{3em}
\setlength{\leftmargini}{2.15em}
\setlength{\leftmarginii}{2.65em}
\setlength{\leftmarginiii}{3em}
\renewcommand{\labelenumi}{\textbf{\theenumi.}}
\pagestyle{empty}
\raggedbottom
\allowdisplaybreaks
\newenvironment{examproblems}[1]{%
  \begin{enumerate}%
  \setcounter{enumi}{#1}%
  \setlength{\topsep}{0.35em}%
  \setlength{\partopsep}{0pt}%
  \setlength{\itemsep}{0.48em}%
  \setlength{\parsep}{0.18em}%
}{\end{enumerate}}
\newenvironment{examsubparts}{%
  \begin{enumerate}%
  \setlength{\topsep}{0.18em}%
  \setlength{\partopsep}{0pt}%
  \setlength{\itemsep}{0.16em}%
  \setlength{\parsep}{0.1em}%
}{\end{enumerate}}
\newenvironment{examinstructions}{%
  \par\begingroup\itshape
}{%
  \par\endgroup\vspace{0.18em}
}
\makeatletter
\renewcommand\section{\@startsection{section}{1}{0pt}%
  {-1.25ex plus -.25ex minus -.1ex}%
  {0.55ex plus .1ex}%
  {\normalfont\large\bfseries}}
\renewcommand{\@maketitle}{%
  \newpage\null\vskip -1.2em%
  {\footnotesize\bfseries
    \href{https://www.ufl.edu/}{University of Florida}, \href{https://math.ufl.edu/}{Department of Mathematics}\par}%
  \vskip 0.18em%
  \tagstructbegin{tag=Title}%
    {\LARGE\bfseries\href{https://gma.math.ufl.edu/exams/combinatorics-phd/}{Combinatorics PhD Exam}, May 1995\par}%
  \tagstructend%
  \vskip 0.35em%
  \tagmcbegin{artifact}%
    \hrule height 1pt%
  \tagmcend%
  \vskip 0.55em%
}
\makeatother
\title{Combinatorics PhD Exam, May 1995}
\author{}
\date{}
\begin{document}
\maketitle
\thispagestyle{empty}
\begin{examproblems}{0}
\item
This problem has two parts.
\begin{examsubparts}
\item[\textnormal{(a)}]
Let \(A_1,A_2,\ldots,A_k\) be subsets of \([n]\) such that \(A_i\cap A_j\ne\varnothing\), for all \(i,j\). Show that 
\[
k\le 2^{n-1}.
\]
 Is this inequality best possible?
\item[\textnormal{(b)}]
How many pairs \((A,B)\) of subsets of \([n]\) are there such that \(A\cap B=\varnothing\)?
\end{examsubparts}
\item
Use inclusion–exclusion to determine the number of monic polynomials of degree~\(n\) with no roots in \(\mathbb{Z}_p[x]\).
\item
This problem has two parts.
\begin{examsubparts}
\item[\textnormal{(a)}]
Prove that the number of partitions of~\(n\) into at most~\(k\) parts equals the number of partitions of~\(n\) into parts of size at most~\(k\).
\item[\textnormal{(b)}]
Prove that the number of partitions of~\(n\) into odd parts is equal to the number of partitions of~\(n\) into unequal parts. (Hint: use generating functions).
\end{examsubparts}
\item
Prove that the number of~\(k\)-dimensional subspaces of the~\(n\) dimensional vector space over \(\mathbb{F}_q\) is 
\[
\binom{n}{k}_q=\frac{(q^n-1)\cdots(q^{n-k+1}-1)}{(q^k-1)\cdots(q-1)}.
\]
\item
Prove that \(K_{33}\) is not planar.
\item
Prove that every regular bipartite graph has a perfect matching. Show that this is not true in general for regular graphs of even order.
\item
Prove that in a binary self dual code, either all the vectors have weight divisible by 4 or half have even weight not divisible by 4 and half have weight divisible by 4.
\item
Let~\(H\) be a parity check matrix for a binary \((n,k,d)\)-code~\(C\) with \(n\ge 4\). If the columns of~\(H\) are distinct and all have odd weight, then \(d\ge 4\).
\item
This problem has four parts.
\begin{examsubparts}
\item[\textnormal{(a)}]
Show the nonexistence of a difference set with parameters \((31,10,3)\).
\item[\textnormal{(b)}]
Show the nonexistence of a symmetric design with parameters \((29,8,2)\). (Hint: consider the appropriate equation modulo 3)
\item[\textnormal{(c)}]
Show the nonexistence of a Steiner system \(S(3,6,11)\).
\item[\textnormal{(d)}]
Find a \((13,4,1)\) difference set~\(D\). The set of translates of~\(D\) form the blocks of a design. Find the parameters \(b,v,k,r,\lambda\).
\end{examsubparts}
\item
Let~\(A\) be an affine plane of order~\(n\). Write \(B_1\sim B_2\) if lines \(B_1\) and \(B_2\) are the same or have no points in common. Show that~\(\sim\) is an equivalence relation. Show that there exists a projective plane~\(\pi\) such that~\(A\) is obtained from~\(\pi\) by deleting one line and all the points on that line.
\end{examproblems}
\end{document}
